\chapter{What Limits the Kernel}\label{r03:ch:root-cause}

\emph{Four measurements locate the limit. It is not memory bandwidth but memory latency: the
X60's hardware prefetcher follows about 8 streams, and the kernel has 42. The arithmetic limit
of the vector unit is \sx{8} over scalar in f32 and \sx{4} in f64.}

All measurements in this chapter use the program as it was before this work: the RVV
variant \code{explicit_cse} and the scalar variant \code{scalar_cse} (\cref{r03:tab:variants}),
pull streaming, compiled with GCC~14, f32, on core~2.

\section{Not bandwidth}\label{r03:sec:bw}

If the kernel were limited by memory bandwidth, it would move data close to the rate the
memory system can deliver, and huge pages, which reduce address-translation overhead, would
matter little. Three observations contradict that:
\begin{itemize}
  \item The throughput hardly depends on the grid size: \code{explicit_cse} gives 2.09, 2.18
        and 2.34~MLUP/s at $16^3$, $32^3$ and $64^3$ cells. All three grids exceed the
        512~KiB L2 cache, so this alone does not rule out memory as the limit, but it does
        not look like a bandwidth limit either.
  \item At $64^3$ the vector kernel spends about 685 cycles per cell at an IPC of 0.13, with
        only about 15 L1 misses per iteration of 8 cells. It moves about 0.4~GB/s. A
        two-dimensional LBM kernel (D2Q9, 9 PDFs per cell) moved 1.6~GB/s on one core of the
        same board in the earlier work.
  \item Huge pages help little: with \code{GLIBC_TUNABLES=glibc.malloc.hugetlb=1}, which
        makes the C library request transparent huge pages, the vector kernel gains 7\,\% and
        the scalar kernel 8\,\%, with half as many TLB misses.
\end{itemize}
The bandwidth explanation was therefore withdrawn for the single-core case.

\section{Where the cycles go}\label{r03:sec:cycles}

\paragraph{Profile.} \code{perf record -e cpu-clock} at $64^3$: the kernel takes 88.9\,\% of
the samples for the vector variant and 91.7\,\% for the scalar one. The halo exchange
(\code{localCopy} of the pull pack info) takes 5.4\,\% and 4.1\,\%. The hot loop has 371
instructions per 8 cells. Among them is a single division (\code{vfdiv}): the other nine
divisions of the source are by constants, and the compiler moves them out of the loop. Most
samples sit on the 19 PDF loads and on the additions that use the loaded values.

\paragraph{Hardware counters.} The X60's vendor events, counted with
\ghws{misc/x60-lab/pmu_stat.sh}. The start-up of the program is subtracted: each event is also
counted in a run with zero time steps (\ghws{misc/x60-lab/pmu_report.py}).

\begin{table}[htbp]
  \centering\small
  \caption{Hardware counters per cell update, benchmark application, $64^3$ cells, f32, GCC.}\label{r03:tab:pmu1}
  \begin{tabular}{@{}lrr@{}}
    \toprule
    \textbf{Counter (meaning)} & \textbf{explicit\_cse} & \textbf{scalar\_cse}\\
    \midrule
    cycles & 682 & 1033\\
    instructions & 61.6 & 337\\
    \code{stalled_cycle_backend} (backend stall cycles) & 630 & 770\\
    \code{eu_lsu_load_full} (cycles with the load queue full) & 420 & 297\\
    \code{l1d_miss} / \code{l2_miss} (misses in L1D / in L2) & 3.49 / 2.97 & 3.44 / 2.98\\
    \code{l1d_prefetch_hit} (L1D hits on prefetched lines) & 0.36 & 0.36\\
    \bottomrule
  \end{tabular}
\end{table}

92\,\% of the vector kernel's cycles are backend stalls. Almost every L1 miss is also an L2
miss, so the data comes from DRAM, and the L1 prefetcher is nearly idle. $682 / 2.97 \approx
230$ cycles per DRAM line is close to one full memory latency per line: the misses hardly
overlap. The scalar kernel has the same miss profile. Both kernels are limited by memory
latency; the scalar one also executes 5.5 times as many instructions.

\section{What the X60 memory system does}\label{r03:sec:mem}

\paragraph{Latency.} A random pointer chase over 128~MB (each load's address comes from the
previous load, so the misses cannot overlap) takes 357 cycles per miss with huge pages and
about 500 with 4~KiB pages. With several independent chains, up to about 6 misses overlap;
this caps random access at about 61 cycles per line.

\paragraph{Streams.} The probe runs $R$ read and $W$ write streams at the same time, each over
4~MB of f32 vectors, with huge pages. The start addresses are offset by 4160~bytes from each
other, so that the streams do not compete for the same cache sets.

\begin{table}[htbp]
  \centering\small
  \caption{Cost of concurrent streams on one X60 core (log \texttt{p2-stream-sweep.txt}).}\label{r03:tab:streams}
  \begin{tabular}{@{}lrr@{}}
    \toprule
    \textbf{Streams} & \textbf{cycles per 64-byte line} & \textbf{GB/s}\\
    \midrule
    1 read & 45 & 2.2\\
    8 read & \textbf{17} & \textbf{6.0}\\
    19 read & 75 & 1.4\\
    24 / 32 read & 364 & 0.28\\
    8 read + 8 write & 37 & 2.8\\
    19 read + 23 write (the pattern of the pull kernel) & \textbf{181--184} & \textbf{0.56}\\
    \bottomrule
  \end{tabular}
\end{table}

The hardware prefetcher handles about 8 read streams well and gives up beyond about 20; then
every line costs a full DRAM latency. The pull kernel, with 19 read and 23 write streams, is
exactly in that regime. Scalar loads collapse at the same point, so the effect is not specific
to vector loads.

\begin{figure}[htbp]
  \centering
  % Cost of concurrent unit-stride read streams on one X60 core (board-logs/p2-stream-sweep.txt
% and p2-prefetch.txt; mem stream R W, f32 vle32, 4 MB per stream, THP, bases staggered by 4160 B).
\begin{tikzpicture}
  \begin{semilogyaxis}[
      width=0.92\linewidth, height=6.2cm,
      xlabel={concurrent read streams $R$ (no writes)},
      ylabel={cycles per 64-byte line},
      xmin=0, xmax=34, ymin=10, ymax=600,
      xtick={1,2,4,8,12,16,19,24,32},
      ytick={10,20,50,100,200,500},
      yticklabels={10,20,50,100,200,500},
      grid=major, grid style={psgray!20},
      tick label style={font=\scriptsize}, label style={font=\footnotesize},
      every axis plot/.append style={thick},
    ]
    % the hardware prefetcher's good zone
    \fill[r03landed!8] (axis cs:0.5,10) rectangle (axis cs:8.5,600);
    \node[r03label,anchor=north] at (axis cs:4.5,560) {prefetcher\\follows};
    \fill[r03rejected!6] (axis cs:19.5,10) rectangle (axis cs:34,600);
    \node[r03label,anchor=south] at (axis cs:27,14) {prefetcher gives up:\\full DRAM latency per line};
    \addplot[color=psblue,mark=*,mark size=1.6pt] coordinates {
      (1,45.1) (2,29.5) (4,25.0) (8,16.9) (12,36.7) (16,34.7) (19,75.3) (24,363.6) (32,363.7)};
    \node[r03label,anchor=west] at (axis cs:8.3,15.5) {8 streams: 17 (6.0 GB/s)};
    % the pull sweep's 19 read + 23 write streams, without and with software prefetch
    \draw[r03rejected,thick,dashed] (axis cs:0,184.3) -- (axis cs:34,184.3);
    \node[r03label,anchor=south west,text=r03rejected!80!black] at (axis cs:0.4,184.3)
      {19 read + 23 write (the pull sweep): 184};
    \draw[r03landed,thick,dashed] (axis cs:0,90.4) -- (axis cs:34,90.4);
    \node[r03label,anchor=south west,text=r03landed!70!black] at (axis cs:9.5,90.4)
      {same, with \texttt{prefetch.r} 128~B ahead: 90};
  \end{semilogyaxis}
\end{tikzpicture}

  \caption[Cost of concurrent streams on the X60]{Cycles per 64-byte line against the number
    of concurrent read streams (logs \texttt{p2-stream-sweep.txt} and
    \texttt{p2-prefetch.txt}). Dashed lines: the 42-stream pattern of the pull kernel, without
    and with software prefetch.}\label{r03:fig:streams}
\end{figure}

\paragraph{Remedies tried on the 19 + 23 pattern.}
\begin{itemize}
  \item Software prefetch (\code{prefetch.r}) 128~bytes ahead of each read stream: from 181 to
        90 cycles per line. \code{prefetch.r} and \code{prefetch.w} 256~bytes ahead: 86.
  \item \textbf{The X60 executes \code{prefetch.r} and \code{prefetch.w} although its ISA
        string does not list Zicbop.} Because they are HINT encodings
        (\cref{r03:sec:bg-rvv}), they are safe on any RISC-V core.
  \item \code{cbo.zero} on the destination lines, which avoids the write-allocate read, does
        not help: 162 alone, 228 together with prefetch. Loading a value from each line ahead
        of time (``touch-ahead'') makes it worse: 195--211. Turning transparent huge pages off
        changes nothing: 179.
\end{itemize}

\begin{keyidea}[Root cause]
The X60's hardware prefetcher fails on 42 concurrent streams, and the in-order core keeps only
about 6 misses in flight. Two kinds of optimisation follow:
\begin{enumerate}[label=(\alph*)]
  \item software prefetch, which a code generator can emit cheaply;
  \item fewer concurrent streams: in-place streaming, or splitting the update into passes that
        touch only a few streams each.
\end{enumerate}
\end{keyidea}

\section{Instruction costs and the arithmetic limit}\label{r03:sec:ops}

\begin{table}[htbp]
  \centering\small
  \caption{X60 cycles per instruction, LMUL~1, VLEN~256 (log \texttt{p2-ops.txt}). Throughput:
    cycles between two independent instructions; latency: cycles until the result can be
    used.}\label{r03:tab:ops}
  \begin{tabularx}{\linewidth}{@{}Xrrl@{}}
    \toprule
    \textbf{Instruction} & \textbf{e32} & \textbf{e64} & \textbf{scalar}\\
    \midrule
    \code{vfadd}/\code{vfmul}/\code{vfmacc} throughput & 1.0 & 1.0 & 1.0 (f32), 1.0/1.3 (f64)\\
    latency add / mul / FMA & 4 / 4 / 5 & 4 / 5 / 6 & 4 (\code{fadd.s})\\
    .vf forms, throughput & 1.0 & 1.0 &\\
    \code{vfdiv} throughput / latency & 38 / 38 & 40 / 12 & 15 (\code{fdiv.s}), 24 (\code{fdiv.d})\\
    \code{vfmv.v.f} (splat) & 4 & 4 &\\
    \code{vle} from L1, aligned / not aligned to the element size & 2 / 4 & 2 / 4 & 1.3 (\code{flw})\\
    \code{vse} to L1 & 2.6 & 2.6 & 1.5 (\code{fsw})\\
    \bottomrule
  \end{tabularx}
\end{table}

The vector unit completes one full-width (256-bit) floating-point instruction per cycle, at
e32 and at e64. That is 8 f32 or 4 f64 operations per cycle, against one for the scalar unit.
The arithmetic limit of vectorization is therefore \sx{8} in single and \sx{4} in double
precision: the goals of this work are the hardware's limits. Vector memory instructions are
relatively expensive (one 8-lane load costs as much as 2 to 4 scalar loads), and the one
division per 8 cells costs 38 cycles.

Two errors in the probe itself were found and fixed while measuring. Linux discards the vector
configuration on every system call, so the system call that starts the cycle counter must come
before the instruction that sets up the vector unit (\code{vsetvli}); otherwise the next vector
instruction aborts the program with \code{SIGILL}. And 64-bit vector loads from an address that
is only 4-byte aligned abort it with \code{SIGBUS}. A third error, in the latency probe of
\cref{r03:sec:mem}: its first version kept the chain pointers in an array, which the compiler
placed on the stack. The store back to the stack serialised the chains, so that misses seemed
not to overlap at all from 4 chains on. The fixed version keeps the pointers in registers.
